% badness off
% ---=== Para matemáticas ===---
\usepackage{amsmath}                        % === Para usar matemáticas
\usepackage{cancel}                         % === Para tachar en entornos matemáticos
\usepackage{amsfonts}                       % === Para texto matemático
\usepackage{amssymb}                        % === Para más símbolos matemáticos
\usepackage{mathrsfs}                       % === Para símbolos matemáticos
\usepackage{derivative}                     % === Para comandos de derivadas


% ---=== Para el formato del texto ===---
\usepackage{ragged2e}                       % === Para elegir la posición del texto
\usepackage[absolute,overlay]{textpos}      % === Para crear bloques de texto
\usepackage[
  letterpaper,
  top=1.5in,
  bottom=1in,
  left=1in,
  right=1in,
  headheight=1in,
  headsep=8mm,
  footskip=12mm
]{geometry}                       % === Para el formato de las páginas
\usepackage{listings}                       % === Para el formato de códigos
\usepackage[spanish,es-tabla]{babel}        % === Para el idioma español
\usepackage{float}                          % === Para forzar la posición de entornos gráficos
\usepackage{hyperref}                       % === Para hiper-vínculos
\usepackage{bookmark}                       % === Para complementar hyperref
\bookmarksetup{open,numbered}
\usepackage{caption}                        % === Para agregar captions
\usepackage[T1]{fontenc}                    % === Para la codificación de la fuente
\usepackage[utf8]{inputenc}                 % === Para los símbolos del español
\usepackage{enumitem}                       % === Para hacer listas
\usepackage{fancyhdr}                       % === Para personalizar el encabezado
\usepackage{lmodern}                        % === Para fuentes tipográficas
\usepackage{setspace}                       % === Para espaciar lineas
\usepackage{wasysym}                        % === Para símbolos astronómicos y de otros tipos
% \usepackage{titlesec}                       % === Para dar formato a las secciones
% \usepackage{soul}                           % === Para tachar texto
% \usepackage{titling}                        % === Para dar formato automático a los titulos


% ---=== Para entornos gráficos ===---

\usepackage{tcolorbox}                      % === Para crear cajas de colores
\usepackage{graphicx}                       % === Para añadir imágenes
\usepackage{pgfplots}                       % === Para graficar
\pgfplotsset{compat=1.15}
\usepackage{tikz}                           % === Para dibujar
\usetikzlibrary{arrows, babel, positioning, fadings, calc}
\usepackage{array}                          % === Para complementar tabular
\usepackage[dvipsnames]{xcolor}             % === Para tener más colores
\usepackage{longtable}                      % === Para tablas de más de una página
\usepackage{tabularx}                       % === Para crear tablas
\usepackage{colortbl}                       % === Para colorear tablas
% \usepackage{rotating}                       % === Para rotar entornos gráficos


\decimalpoint{}                             % Para usar "." como punto decimal
\setlength{\parindent}{0cm}                 % Para iniciar los párrafos sin indentación

\tcbuselibrary{listings, skins, breakable}  % Para... algo... de listins... creo...
